\section{Metodologia}

Para responder se a liquidez recente de um clube comprador eleva causalmente o sobrepreço que ele
paga, adotamos um \textit{pipeline} em duas etapas encadeadas. A primeira etapa estima, por
aprendizado de máquina, o \textit{preço justo} de cada transferência a partir de atributos do
jogador, isolando como resíduo o sobrepreço pago. A segunda etapa toma esse resíduo como desfecho e
estima o efeito causal da liquidez do comprador sobre ele, neutralizando confundidores por meio de
\textit{Double Machine Learning}. Esta seção descreve os dados, o pré-processamento e as duas
etapas de modelagem.

\subsection{Conjunto de dados e pré-processamento}

Os dados foram coletados do \textit{Transfermarkt} por meio de um \textit{scraper} próprio, cobrindo
\textbf{oito temporadas} (de 2017 a 2025, com exceção de 2020, não coletada) das \textbf{sete
principais ligas europeias} (Inglaterra, Espanha, Itália, Alemanha, França, Portugal e Bélgica).
A base bruta contém \textbf{44.627 movimentações}, incluindo empréstimos, transferências gratuitas e
fins de empréstimo.

\textbf{Filtragem.} Como o objeto de estudo é o preço negociado, retemos apenas \textit{compras
pagas} com \textit{fee} e valor de mercado válidos. Transferências abaixo de \textbf{€250 mil} são
descartadas: nessa faixa o prêmio relativo é fortemente negativo e de baixa variância, indicando
ausência de dinâmica real de mercado (movimentações administrativas). Após a filtragem e a remoção
de registros incompletos, o conjunto de modelagem contém \textbf{5.146 transferências}.

\textbf{Transformações.} Os valores monetários têm forte assimetria à direita; aplicamos a
transformação logarítmica ($\log(1+x)$) ao \textit{fee} e ao valor de mercado para estabilizar a
escala. O prêmio bruto é \textit{winsorizado} nos percentis 1 e 99 para conter caudas extremas
(que chegam a milhares de pontos percentuais).

\textbf{Variáveis derivadas do clube.} Para cada clube comprador em cada temporada, agregamos a
\textit{receita de vendas} (soma dos \textit{fees} recebidos), o volume de compras (gasto total e
número de contratações) e o saldo líquido. A receita de vendas é a variável de tratamento; as
demais permitem separar o efeito de \textit{ter caixa} do comportamento habitual de gasto do clube.

\textbf{Variáveis de rede.} O mercado é modelado como um grafo dirigido e ponderado, reconstruído
por temporada, em que os nós são clubes, as arestas representam transferências (vendedor para
comprador) e o peso é o \textit{fee}. De cada clube extraímos métricas de centralidade — em
particular o \textit{PageRank} e os graus de entrada e saída — que capturam sua posição estrutural
no ecossistema e servem como controles.

\subsection{Etapa 1: modelo hedônico do preço de referência}

A primeira etapa implementa um \textit{modelo hedônico}, que decompõe o preço de um bem na soma do
valor de seus atributos. Aqui, o preço de uma transferência é explicado por características do
jogador e do contexto de mercado. A variável-alvo é o logaritmo do \textit{fee}, e o que o modelo
\textbf{não} consegue explicar — o resíduo — é interpretado como o \textit{sobrepreço de
reinvestimento}:
%
\begin{equation}
Y_{i} \;=\; \ln(P_i) \;-\; \ln(\hat{P}_i),
\label{eq:residuo}
\end{equation}
%
onde $P_i$ é o \textit{fee} pago na transferência $i$ e $\hat{P}_i$ é o preço previsto pelo modelo.
Por construção, $Y_i$ está livre das características intrínsecas do atleta; o que resta são forças
conjunturais da negociação — exatamente o objeto da Etapa 2.

\textbf{Covariáveis.} O modelo usa \textbf{19 variáveis}, todas descrevendo o jogador ou o contexto
estrutural de mercado, \textbf{nunca} o comportamento financeiro do comprador: idade e idade ao
quadrado (capturando a não-linearidade do ciclo de carreira), o logaritmo do valor de mercado,
indicadores de posição, seis \textit{dummies} de liga e sete \textit{dummies} de temporada (efeitos
fixos que absorvem a inflação de mercado). Essa separação é deliberada: reservar o comportamento
financeiro do clube para a etapa causal garante que o resíduo não o contenha de antemão.

\textbf{Modelos e avaliação.} Comparamos quatro algoritmos de regressão — \textit{Random Forest}
\cite{breiman2001}, \textit{XGBoost}, \textit{LightGBM} e \textit{Support Vector Regression} —
adotando um \textit{split} temporal: treino nas temporadas de 2017 a 2024 (4.396 observações) e
teste \textit{holdout} em 2025 (750 observações). Esse desenho simula o uso real — treinar no
passado e prever o futuro — e oferece um teste de generalização mais honesto que a partição
aleatória. A métrica principal é o RMSE, que penaliza erros grandes, relevante dada a amplitude dos
\textit{fees}. Como \textit{baseline}, avaliamos a predição usando apenas o valor de mercado, sem
aprendizado adicional. A importância das variáveis é interpretada com valores SHAP
\cite{lundberg2017}.

\textbf{Prevenção de vazamento.} O resíduo $Y_i$ que alimenta a Etapa 2 é calculado de forma
\textit{out-of-fold} (via \textit{cross-validation}): a previsão de cada observação provém de um
modelo treinado sem ela. Isso evita o vazamento que ocorreria ao usar previsões \textit{in-sample},
o qual comprimiria artificialmente a variância do resíduo e contaminaria a etapa causal.

\subsection{Etapa 2: estimação causal por Double Machine Learning}

Na segunda etapa, estimamos o efeito causal da liquidez do comprador sobre o sobrepreço. O
\textbf{tratamento} $D$ é o logaritmo da receita de vendas do clube na temporada; o \textbf{desfecho}
$Y$ é o resíduo hedônico da Etapa 1. O desafio é que diversos fatores estruturais afetam tanto a
receita quanto o sobrepreço — sobretudo o fato de que clubes grandes vendem caro \textit{e} compram
caro. Para neutralizá-los, empregamos \textit{Double/Debiased Machine Learning} (DML)
\cite{chernozhukov2018}, que assume o modelo parcialmente linear
%
\begin{equation}
Y \;=\; \theta_0\, D \;+\; g(W) \;+\; U, \qquad D \;=\; m(W) \;+\; V,
\label{eq:dml}
\end{equation}
%
em que $W$ é o vetor de confundidores e $\theta_0$ é o parâmetro causal de interesse. A ideia central
é remover de $Y$ e de $D$ a parte explicada por $W$ — usando modelos flexíveis de aprendizado de
máquina para estimar $g$ e $m$ — e estimar $\theta_0$ a partir dos resíduos ortogonalizados. Para
evitar viés de sobreajuste, as funções $g$ e $m$ são ajustadas por \textit{cross-fitting} em cinco
partições, com previsões sempre feitas fora da partição de treino. Os erros-padrão são
\textbf{clusterizados por clube e temporada}, refletindo que o tratamento varia no nível
clube-temporada.

\textbf{Confundidores.} O vetor $W$ reúne \textbf{19 variáveis}: as métricas de rede (\textit{PageRank}
e graus), o porte do elenco (tamanho, idade média, número de jogadores de seleção), além das
\textit{dummies} de liga e de temporada. Esses controles capturam o porte e a posição estrutural do
clube. Cabe registrar uma limitação: o vetor $W$ não inclui o volume financeiro \textit{direto}
(gasto total, número de compras); o confundidor de ``clube rico'' é, portanto, controlado por
\textit{proxies} estruturais, e não pelo gasto em si — ponto retomado na análise de robustez.

\textbf{Heterogeneidade e índice de vulnerabilidade.} Além do efeito médio, estimamos como o efeito
varia entre clubes. Para isso, usamos o \textit{R-learner} \cite{nie2021}, que produz um efeito
causal condicional (CATE) por observação. Agregando os efeitos por clube comprador, definimos o
\textit{Índice de Vulnerabilidade de Barganha} (IVB), que posiciona cada clube numa escala de 0 a 1:
valores altos identificam ``presas fáceis'' (clubes que pagam mais sobrepreço quando capitalizados)
e valores baixos, ``negociadores disciplinados''.

\subsection{Protocolo de validação e robustez}

Como o efeito estimado é pequeno e a identificação repousa em suposições, submetemos os resultados a
uma bateria de testes, cujos resultados são reportados na Seção~\ref{sec:resultados}. \textbf{Testes
placebo} substituem o tratamento por ruído (permutação e ruído gaussiano), verificando que o modelo
não detecta efeito espúrio. Um \textbf{teste de tratamento defasado} usa a receita da temporada
anterior — predeterminada em relação à compra atual — para avaliar causalidade reversa. Os efeitos
por temporada são submetidos a \textbf{correção de múltiplas comparações} (Bonferroni e
\textit{false discovery rate} \cite{benjamini1995}). Um \textbf{placebo de permutação intra-safra}
testa a robustez dos efeitos sazonais. Avaliamos ainda a \textbf{sensibilidade ao limiar de
\textit{fee}}, reconstruindo o \textit{pipeline} para diferentes cortes, e quantificamos a
fragilidade a confundidores omitidos pelo \textbf{\textit{Robustness Value}} \cite{cinelli2020}.
Por fim, um \textbf{teste de \textit{over-control}} reintroduz o volume financeiro no vetor de
confundidores, e \textbf{especificações alternativas do tratamento} (número de vendas e indicador de
venda de grande porte) examinam o mecanismo subjacente.
